% Table 5 (3.5 Survival stratification)
% 자동 생성: tools/make_paper_tables.py — 수치를 여기서 고치지 말고 스크립트를 다시 돌릴 것.
% 단독 문서다(캡처용). Overleaf 새 프로젝트에서 pdfLaTeX 로 컴파일한다.
% 기존 main.tex 에 넣을 때는 이 파일이 아니라 overleaf_tables.tex 를 쓴다.
\documentclass[border=8pt]{standalone}
\usepackage[T1]{fontenc}
\usepackage{newtxtext,newtxmath}
\usepackage{booktabs,multirow,array,pifont}
% --- 캡처할 때 제목·주석을 빼려면 아래 두 줄을 \...false 로 바꾼다 ---
\newif\ifshowtitle \showtitletrue
\newif\ifshownotes \shownotestrue
\newsavebox\tablebox
\begin{document}
\sbox\tablebox{%
\begin{tabular}{@{}ll cc c cc c c@{}}
\toprule
Endpoint & Risk group & Patients & Events & Median, months & 1-year & 2-year & HR (95\% CI) & \textit{P} \\
 & & & & (95\% CI) & rate, \% & rate, \% & & \\
\midrule
\multirow{2}{*}{OS} & Low risk & 112 & 83 & 18.9 (15.5--24.3) & 67.9 & 42.7 & 1 (reference) & \multirow{2}{*}{$<$ 0.001} \\
 & High risk & 126 & 116 & 7.7 (6.4--9.6) & 29.4 & 12.7 & 2.32 (1.74--3.10) &  \\
\addlinespace
\multirow{2}{*}{PFS} & Low risk & 122 & 101 & 8.5 (6.6--10.1) & 34.4 & 20.2 & 1 (reference) & \multirow{2}{*}{$<$ 0.001} \\
 & High risk & 116 & 111 & 5.2 (4.7--5.6) & 10.3 & 5.9 & 2.03 (1.54--2.67) &  \\
\bottomrule
\end{tabular}%
}%
\begin{minipage}{\wd\tablebox}
\ifshowtitle\noindent \textbf{Table 5.} Survival According to the Risk Group Defined by the Proposed Model\par\vspace{4pt}\fi
\noindent\usebox\tablebox
\ifshownotes\par\vspace{3pt}\footnotesize\noindent High risk, risk score above the median of the training patients of each fold; the five test folds were pooled (not an independent test set). \textit{P}, log-rank test; HR, Cox model. CI = confidence interval, HR = hazard ratio, OS = overall survival, PFS = progression-free survival\par\fi
\end{minipage}
\end{document}
